\section{Los signos}\label{sec:signos}

Hay veintiséis signos activos, repartidos en cuatro familias. Cada familia tiene
un trabajo distinto en la línea, así que al mirar un signo ya sé dónde debería
estar y qué hace.

\subsection{Conectores: cómo se relacionan dos ideas}
Van \emph{entre} dos ideas y se leen de izquierda a derecha. Los que llevan
punta de flecha tienen dirección: cambiar el orden de las ideas cambia lo que
se dice. Los que no llevan punta son simétricos.\par\vspace{2mm}
\KalaTarjetas{K23}{K21}{K24}{K03}
\KalaTarjetas{K26}{K27}{K25}{K28}
\noindent\KalaTarjeta{K29}\hfill\KalaTarjetaVacia\hfill\KalaTarjetaVacia\hfill\KalaTarjetaVacia\par\vspace{3mm}

\subsection{Categorías: qué tipo de cosa es un nombre}
Se escriben pequeñas \emph{delante} de un nombre, y solo cuando ayudan: la
primera vez que aparece o cuando podría confundirse. «\kss{K30} Turing» deja
claro que hablo de la persona y no de la máquina.\par\vspace{2mm}
\KalaTarjetas{K30}{K08}{K07}{K04}
\KalaTarjetas{K01}{K02}{K05}{K06}

\subsection{Marcadores: cómo leer una afirmación}
Dicen de dónde sale una idea o qué hacer con ella. Se escriben dentro de la
línea, delante de lo que marcan.\par\vspace{2mm}
\KalaTarjetas{K09}{K16}{K19}{K10}

\subsection{Margen: marcas para repasar}
Van en la columna del margen y afectan a toda la línea. Al repasar, recorro el
margen y encuentro enseguida lo importante, mis ideas y lo pendiente.\par\vspace{2mm}
\KalaTarjetas{K31}{K35}{K32}{K33}
\noindent\KalaTarjeta{K34}\hfill\KalaTarjetaVacia\hfill\KalaTarjetaVacia\hfill\KalaTarjetaVacia\par\vspace{3mm}

{\small\color{KalaSuave}Pregunta \kss{K19} puede ir en el margen, para dudar de
la línea entera, o dentro de la línea, para dudar solo de una parte.}

\clearpage
\subsection{Catálogo}\label{sec:catalogo}
\begin{longtable}{@{}p{29mm}>{\raggedright\arraybackslash}p{\dimexpr\linewidth-29mm-2\tabcolsep\relax}@{}}
\toprule \textbf{Signo} & \textbf{Qué significa} \\ \midrule \endfirsthead
\toprule \textbf{Signo} & \textbf{Qué significa} \\ \midrule \endhead
\bottomrule \endfoot
\KalaCatalogoFila{K23}
\KalaCatalogoFila{K21}
\KalaCatalogoFila{K24}
\KalaCatalogoFila{K03}
\KalaCatalogoFila{K26}
\KalaCatalogoFila{K27}
\KalaCatalogoFila{K25}
\KalaCatalogoFila{K28}
\KalaCatalogoFila{K29}
\KalaCatalogoFila{K30}
\KalaCatalogoFila{K08}
\KalaCatalogoFila{K07}
\KalaCatalogoFila{K04}
\KalaCatalogoFila{K01}
\KalaCatalogoFila{K02}
\KalaCatalogoFila{K05}
\KalaCatalogoFila{K06}
\KalaCatalogoFila{K09}
\KalaCatalogoFila{K16}
\KalaCatalogoFila{K19}
\KalaCatalogoFila{K10}
\KalaCatalogoFila{K31}
\KalaCatalogoFila{K35}
\KalaCatalogoFila{K32}
\KalaCatalogoFila{K33}
\KalaCatalogoFila{K34}
\end{longtable}
Los códigos no son consecutivos porque conservan los de la versión anterior:
un código nunca se reutiliza con otro significado. Los que faltan están en el
anexo~\ref{sec:retirados}.

\clearpage
\subsection{Cómo se trazan}\label{sec:trazado}
Cada número es un trazo sin levantar el lápiz. Izquierda, derecha, arriba y
abajo se refieren al signo tal como aparece impreso. El orden es una sugerencia:
puedo cambiarlo si el signo sigue reconociéndose.

Tamaño: el alto de una letra mayúscula de mi letra, más o menos. Dejo un poco
de espacio antes y después de cada conector para que no se pegue a las
palabras vecinas.

% Orden orientativo de ejecución; una fila por signo canónico.
\newcommand{\KalaTrazoFila}[2]{%
\KalaSimbolo[13mm]{#1}\quad\texttt{#1} &
\textbf{\csname KalaNombre#1\endcsname}
(\csname KalaTrayectos#1\endcsname\
\ifnum\csname KalaTrayectos#1\endcsname=1 trayecto\else trayectos\fi). #2\\[7pt]}
\begin{longtable}{@{}p{32mm}>{\raggedright\arraybackslash}p{\dimexpr\linewidth-32mm-2\tabcolsep\relax}@{}}
\toprule Modelo & Orden de trazado \\ \midrule\endhead
\bottomrule\endfoot
\KalaTrazoFila{K01}{1. Bajo por la izquierda y curvo hacia la derecha en la base.
2. Dibujo el ángulo superior desde la izquierda hacia la derecha y bajo.}
\KalaTrazoFila{K02}{1. Desde abajo a la izquierda, subo curvando por arriba hacia
la derecha y vuelvo a la izquierda hasta el centro. 2. Bajo desde ese extremo.}
\KalaTrazoFila{K03}{1. Trazo la curva ascendente de izquierda a derecha y termino
horizontal. 2. Dibujo la punta desde arriba a la izquierda, pasando por el
extremo derecho y bajando a la izquierda.}
\KalaTrazoFila{K04}{1. Desde abajo a la izquierda, dibujo una onda con arco
superior, valle y subida derecha. 2. Añado la punta superior de izquierda a derecha.}
\KalaTrazoFila{K05}{1. Bajo por la izquierda y giro a la derecha en la base.
2. Añado la barra superior hacia la derecha. 3. Añado la barra intermedia.}
\KalaTrazoFila{K06}{1. Arco izquierdo abierto, de arriba abajo.
2. Zigzag derecho: derecha, diagonal abajo a la izquierda y derecha.}
\KalaTrazoFila{K07}{1. Desde arriba a la izquierda, curvo por la derecha y vuelvo
a la izquierda por abajo. 2. Trazo la L interior bajando y girando a la derecha.}
\KalaTrazoFila{K08}{1. Cierro el rombo: izquierda, arriba, derecha, abajo,
izquierda. 2. Barra interior hacia la derecha hasta tocar el lado oblicuo,
sin sobresalir.}
\KalaTrazoFila{K09}{1. Corchete izquierdo: derecha, abajo, izquierda.
2. Barra derecha superior. 3. Barra derecha inferior.}
\KalaTrazoFila{K10}{1. Desde abajo a la izquierda, subo, redondeo la esquina,
sigo a la derecha, bajo y vuelvo hacia el interior a la izquierda.}
\KalaTrazoFila{K11}{1. Arco desde arriba a la izquierda por la derecha hasta
abajo. 2. Diagonal desde abajo a la izquierda hacia arriba a la derecha.}
\KalaTrazoFila{K12}{1. V desde arriba a la izquierda al centro inferior y arriba
a la derecha. 2. Bajo una cola desde el vértice inferior.}
\KalaTrazoFila{K13}{1. Desde arriba a la izquierda, voy al centro y vuelvo abajo
a la izquierda. 2. Del centro hacia la derecha y un retorno corto arriba a la izquierda.}
\KalaTrazoFila{K14}{1. U desde arriba a la izquierda, base curva y subida derecha.
2. Barra horizontal entre los lados. 3. Vertical central desde arriba, cruzando la barra.}
\KalaTrazoFila{K15}{1. La misma U. 2. Barra corta interior, separada de ambos
lados. Conservo esa separación y no añado la vertical de K14.}
\KalaTrazoFila{K16}{1. Desde abajo a la izquierda: subo, derecha, esquina
inclinada, bajo y cierro hacia la izquierda. 2. Barra interior hacia la derecha.}
\KalaTrazoFila{K17}{1. Ángulo izquierdo hacia el centro y de vuelta abajo a la
izquierda. 2. Ángulo derecho hacia el centro y de vuelta abajo a la derecha.
Dejo separados los vértices; no formo una X.}
\KalaTrazoFila{K18}{1. Desde arriba a la derecha, voy a la izquierda, bajo con
esquinas curvas y vuelvo a la derecha por abajo. 2. Ángulo interior hacia la
derecha y abajo a la izquierda.}
\KalaTrazoFila{K19}{1. Desde abajo a la izquierda, curvo por arriba hacia la
derecha, vuelvo en diagonal abajo a la izquierda y termino abajo a la derecha.}
\KalaTrazoFila{K20}{1. Barra superior hacia la derecha, diagonal abajo a la
izquierda y bajada vertical. 2. Barra intermedia hacia la derecha.}
\KalaTrazoFila{K21}{1. Escalera: subo, derecha, subo, derecha.
2. Punta de flecha desde arriba a la izquierda, extremo y abajo a la izquierda.}
\KalaTrazoFila{K22}{1. Desde arriba a la derecha, rodeo por arriba, izquierda y
abajo, subo por la derecha y termino con la barra hacia dentro. Todo es un
trayecto continuo; no añado una bajada ni repaso el arco.}
\end{longtable}


\clearpage
\subsection{Rasgos que no deben perderse}
Al escribir rápido, cada signo puede deformarse mientras conserve el rasgo que
lo distingue de su vecino más parecido:
\begin{center}
\begin{tabular}{@{}cp{.62\linewidth}@{}}
\toprule Par & Qué los distingue \\ \midrule
\KalaPar{K23}{K21} & Causa es un zigzag en diagonal; Luego, una escalera de ángulos rectos. \\[3mm]
\KalaPar{K24}{K03} & Se vuelve tiene un bucle y termina horizontal; Relación es una curva sin bucle. \\[3mm]
\KalaPar{K25}{K29} & Pero son dos arcos que se dan la espalda; Opuesto, dos círculos unidos por una línea. \\[3mm]
\KalaPar{K26}{K27} & Ejemplo es una lupa con mango; Parte de, un círculo dentro de otro. \\[3mm]
\KalaPar{K30}{K26} & Persona tiene piernas abiertas bajo el círculo; Ejemplo, un mango en diagonal. \\[3mm]
\KalaPar{K34}{K16} & Revisar es un cuadrado vacío; Según tiene la esquina cortada y una barra. \\[3mm]
\KalaPar{K10}{K02} & Ver tiene lados rectos y se cierra por dentro; Propiedad es curva y baja al centro. \\[3mm]
\KalaPar{K05}{K09} & Cantidad es una L con dos barras; Definición, un corchete seguido de un igual. \\[3mm]
\KalaPar{K28}{K25} & Van juntos son dos rectas paralelas; Pero, dos curvas. \\
\bottomrule
\end{tabular}
\end{center}

Como los mezclo con letras, también cuido tres parecidos: Cantidad \kss{K05} no
debe parecer una E (la barra superior no existe), Lista \kss{K06} lleva el arco
separado de la z, y Ejemplo \kss{K26} lleva el punto que lo separa de una Q.
